\section{ARQUITECTURA DEL SISTEMA}

La arquitectura del sistema sigue un patrón desacoplado en capas, diseñado para superar la brecha computacional entre el procesamiento masivo de datos no supervisados y la necesidad de respuestas interactivas en milisegundos en la interfaz web del usuario.

El flujo integral se estructura en dos fases operativas: la \textbf{Fase Offline (Entrenamiento, Optimización y Precomputación)} y la \textbf{Fase Online (Servicio Web REST e Inferencia Interactiva)}. La Figura \ref{fig:arquitectura_sistema} detalla la articulación de las capas del sistema.

\begin{figure}[htbp]
\centering
\begin{tikzpicture}[
    node distance=0.8cm and 0.4cm,
    layer_box/.style={rectangle, draw=primaryblue, fill=primaryblue!5, rounded corners=3pt, line width=1pt, inner sep=8pt},
    module_box/.style={rectangle, draw=secondaryblue!90, fill=white, rounded corners=2pt, font=\scriptsize, align=center, minimum height=0.9cm, text width=3.4cm, line width=0.7pt},
    artifact_box/.style={rectangle, draw=accentorange!90, fill=accentorange!10, rounded corners=2pt, font=\scriptsize, align=center, minimum height=0.9cm, text width=3.4cm, line width=0.7pt},
    flask_box/.style={rectangle, draw=accentgreen!90, fill=accentgreen!10, rounded corners=2pt, font=\scriptsize, align=center, minimum height=0.9cm, text width=3.4cm, line width=0.7pt},
    arrow/.style={->, >=stealth, thick, color=darkgray}
]

    % Capa 1: Capa de Datos
    \node (d1) [module_box] {Dataset Chicago (CPD)\\($324.5\text{ MB}$, $+800\text{k}$ filas)};
    \node (d2) [module_box, right=0.6cm of d1] {Pipeline de Limpieza\\(Haversine, Snappy Parquet)};
    \node (d3) [artifact_box, right=0.6cm of d2] {Corpus Bolivia (137)\\y \texttt{chicago\_clean.parquet}};
    
    % Enmarcar Capa de Datos
    \begin{scope}[on background layer]
        \node (layer1) [layer_box, fit=(d1) (d2) (d3)] {};
        \node [above=0.1cm of layer1.north west, anchor=west, font=\bfseries\footnotesize, color=primaryblue] {1. Capa de Datos e Ingestión};
    \end{scope}

    % Capa 2: Capa Analítica Offline
    \node (m1) [module_box, below=1.2cm of d1] {Clustering Espacial\\(DBSCAN, HDBSCAN, OPTICS)};
    \node (m2) [module_box, below=1.2cm of d2] {Particional \& NLP\\(K-Means, TF-IDF, LDA)};
    \node (m3) [module_box, below=1.2cm of d3] {Probabilístico \& Anomalías\\(GMM EM $2\sigma$, Autoencoder)};

    % Enmarcar Capa de Modelado
    \begin{scope}[on background layer]
        \node (layer2) [layer_box, fit=(m1) (m2) (m3)] {};
        \node [above=0.1cm of layer2.north west, anchor=west, font=\bfseries\footnotesize, color=primaryblue] {2. Capa de Modelado No Supervisado (Offline)};
    \end{scope}

    % Capa 3: Persistencia y Microservicio Flask
    \node (s1) [artifact_box, below=1.2cm of m1] {Modelos Serializados\\(\texttt{.joblib}, centroides, elipses)};
    \node (s2) [flask_box, below=1.2cm of m2] {Servidor Web Flask\\(\texttt{app.py} y Blueprints)};
    \node (s3) [flask_box, below=1.2cm of m3] {Endpoints API REST\\(\texttt{/api/predict\_risk}, etc.)};

    % Enmarcar Capa de Servicio
    \begin{scope}[on background layer]
        \node (layer3) [layer_box, fit=(s1) (s2) (s3)] {};
        \node [above=0.1cm of layer3.north west, anchor=west, font=\bfseries\footnotesize, color=primaryblue] {3. Capa de Servicio y Persistencia (Backend)};
    \end{scope}

    % Capa 4: Presentación y Visualización WebGL
    \node (p1) [module_box, below=1.2cm of s1] {Dashboard Ejecutivo\\(KPIs, Métricas Codo/Silueta)};
    \node (p2) [module_box, below=1.2cm of s2] {Visor Globo 3D WebGL\\(MapLibre GL / Deck.gl)};
    \node (p3) [module_box, below=1.2cm of s3] {Playground Interactivo\\(Simulador GMM y NLP en vivo)};

    % Enmarcar Capa de Presentación
    \begin{scope}[on background layer]
        \node (layer4) [layer_box, fit=(p1) (p2) (p3)] {};
        \node [above=0.1cm of layer4.north west, anchor=west, font=\bfseries\footnotesize, color=primaryblue] {4. Capa de Presentación y Experiencia de Usuario};
    \end{scope}

    % Conexiones Intercapas
    \draw [arrow] (d1) -- (d2);
    \draw [arrow] (d2) -- (d3);
    \draw [arrow] (d3) -- (m2);
    \draw [arrow] (m1) -- (s1);
    \draw [arrow] (m2) -- (s1);
    \draw [arrow] (m3) -- (s1);
    \draw [arrow] (s1) -- (s2);
    \draw [arrow] (s2) -- (s3);
    \draw [arrow] (s2) -- (p2);
    \draw [arrow] (s3) -- (p3);
    \draw [arrow] (s2) -- (p1);

\end{tikzpicture}
\caption{Diagrama de Arquitectura Modular del Sistema en Cuatro Capas.}
\label{fig:arquitectura_sistema}
\figurenote{Elaboración propia. Articula el desacoplamiento entre el cómputo intensivo offline (Data Lakehouse y entrenamiento serializado en joblib) y el servicio interactivo online (Flask API REST y visualizador 3D acelerado por WebGL).}
\end{figure}

\subsection{Descripción Funcional de las Capas}

\begin{enumerate}
    \item \textbf{Capa de Datos e Ingestión:} Almacena los datos maestros en Parquet con compresión Snappy y mantiene la muestra balanceada para renderizado web reactivo. Administra asimismo el repositorio estructurado de noticias de Bolivia con los 137 incidentes geocodificados y los gazetteers urbanos.
    \item \textbf{Capa Analítica Offline:} Ejecuta los scripts de entrenamiento por lotes (\texttt{modules/algorithms/}), optimiza los hiperparámetros de densidad y reconstrucción neuronal, y genera las matrices precalculadas de estabilidad y benchmark.
    \item \textbf{Capa de Servicio y Persistencia (Flask Master):} Carga los modelos serializados binariamente con \texttt{joblib} al iniciar el servidor. A través de Blueprints modulares (\texttt{modules/routes/}), expone los endpoints JSON RESTful detallados en la Tabla \ref{tab:endpoints_api}.
    \item \textbf{Capa de Presentación (Frontend WebGL 3D):} Renderiza el tablero de mando interactivo utilizando HTML5 semántico, hojas de estilo CSS3 y scripts en JavaScript conectados al motor WebGL (MapLibre GL JS / Deck.gl) para proyectar el Globo 3D con conmutación dinámica de capas algorítmicas.
\end{enumerate}

\begin{table}[htbp]
\centering
\small
\begin{tabularx}{\textwidth}{l l p{4.2cm} X}
\toprule
\textbf{Método} & \textbf{Ruta del Endpoint} & \textbf{Entrada / Parámetros} & \textbf{Función y Respuesta JSON} \\
\midrule
\texttt{GET} & \texttt{/} & Ninguno & Renderiza el dashboard principal interactivo. \\
\texttt{GET} & \texttt{/webgis} & \texttt{ctx} (bolivia / chicago) & Carga la consola del visor esférico Globo 3D WebGL. \\
\texttt{POST} & \texttt{/api/predict\_risk} & Coordenadas $(\text{lat}, \text{lon})$ & Retorna la densidad continua $p(x)$ y componentes GMM. \\
\texttt{POST} & \texttt{/api/classify\_text} & Cadena de texto libre & Infiere el clúster semántico TF-IDF y modus operandi. \\
\texttt{GET} & \texttt{/api/benchmark} & \texttt{distrito\_id} (CPD) & Provee métricas auditadas (DBSCAN, HDBSCAN, OPTICS). \\
\texttt{POST} & \texttt{/api/project\_dirichlet} & Texto o preset criminal & Estima las probabilidades posteriores del simplex LDA. \\
\bottomrule
\end{tabularx}
\caption{Catálogo de endpoints y microservicios REST de la plataforma analítica.}
\label{tab:endpoints_api}
\tablenote{Elaboración propia a partir de la especificación de rutas en \texttt{modules/routes/api\_routes.py}. Todos los endpoints operan con tiempos de respuesta en memoria inferiores a 50 milisegundos.}
\end{table}

